\section{The Simultaneous Trie Descent}
\label{app:trie}

The descent keeps comparisons by recording their entering transitions.
Those labels record which block was read and the state from which it
was read, so that a test on blocks can refer back to a cut inside a
block.

\subsection{Transition labels and linked expansions}

A transition reader uses a deterministic table with designated starts.
The label of the first cut is its coordinate start symbol. At a later
cut, the label is the pair of the previous state and the letter just
read. A label determines the state of the cut. The reader tests the
last label of each coordinate for eligibility. For comparisons it has
two admission relations on labels for each ordered pair of coordinates.
One is used when the target rest is nonempty and the other is used
when it is empty. A hole meets the endpoint tests and has no admitted
comparison.

At a given level, classify exactly the last letters of the nonempty
words as markers. The other letters are ordinary. Splitting after every
marker yields a uniquely linked block list for each coordinate, including
the empty list for an empty word. Write a block as $(\mu,x,\nu)$ with front
marker $\mu$, content $x$, and closing marker $\nu$. The front marker is
either a letter from the marker set or a distinct coordinate start
symbol. The expansion writes $x\nu$ for each block. Lists linked from
the same front marker have expansions in the prefix relation exactly
when the lists are in the prefix relation. To see the reverse
implication, scan the expansions to the first marker. The matching
initial pieces give the same content and closing marker, and the
fronts agreed. Repeat for every block of the shorter expansion.

The block transition starts with the front marker. If it is a letter,
read it from the current state. If it is a coordinate start, reset
to the associated start state. Then read the content. On a linked
list, the state after a block is the original state just before the
closing marker. That marker is the next block's front, so induction
over the list proves the claimed state correspondence. A noninitial
label of the block list consists of the state before a block and
the block itself. From it one can reconstruct the original label
after any prefix of the content, including the one at the front.
One can also compute the label of the closing transition. For an
initial label use the start symbol as its closing label, and give
it no interior cuts.

A proper cut of an expansion belongs to the first block whose closing
marker lies at or after the cut. Its rest is the rest of the content,
then the closing marker, then the expansion of the remaining blocks.
For two such cuts, the prefix test holds exactly when the content
rests and the closing markers agree and the following block lists
satisfy the block prefix test. Indeed, the first marker in each rest
must be at the same position, and the lists after it are linked from
the same front marker. The final cut of an expansion has the closing
label of the final block-list cut.

These correspondences give the derived reader. Its endpoint acceptance
tests the closing label. At a comparison of block cuts with nonempty
target rest it admits if there are cuts inside the two indicated
blocks with equal content rests and closing markers whose original
labels are admitted. At a comparison with the final block cut of the
target it admits by the same test, or if the closing label of the
target is admitted as a final target with some label represented
by the host cut. The labels represented by a host cut are the interior
labels of its block and its closing label. The initial cut of a
block list also represents the original start via its closing label.
Each of these represented labels occurs in the expansion at a cut
of that host. Conversely each cut of the expansion has a label
represented by some block cut. Thus every comparison admitted in an
expansion has a derived test detecting it. Each derived test that
admits when the block rests match also supplies an admitted
comparison in the expansion. This proves the equivalence of holes
of the derived reader on linked lists and holes on their expansions.

The junction monitor has a state for each coordinate start and each
marker, and an absorbing failure state. A block step tests its front
against this state and outputs its closing marker on success. With
the start state of its coordinate, a run avoids failure exactly
when its blocks are linked, by induction on the list. Take the product
with the derived reader. In tests on labels forget the monitor state,
and require at endpoints that it has avoided failure. Its holes are
exactly the linked holes of the derived reader.

\subsection{Depth and word actions at one level}

A chain of nonempty rests consists of words each of which is a proper
prefix of the next. Any such chain is counted by the depth test on
its longest word. Conversely the prefixes in any depth test are a
chain of rests when sorted by length. For a chain of at least two
block rests, the first block is the same in each rest. Its front
cannot be a coordinate start, since that would place the first
block at the start of one specific coordinate in every rest. So the
front is a marker $\mu$. Prepend $\mu$ to the expansion of each rest.
These are rests of the expanded tuple and remain a strict chain
by the prefix equivalence for linked lists. The one-letter rest
$\mu$ occurs at the end of some coordinate by choice of the markers.
Adding it extends the chain by one. Thus if the original depth
is $d\ge 2$, the block depth is at most $d-1$. If the original
depth is one, each list has at most one block. A second block
and the following blocks would form a rest with ordinary front
marker $\mu$, giving the rest $\mu$ followed by that expansion
and the shorter rest $\mu$.

For a list of blocks, the induced state map of the derived reader is
either a word action of the old table or a constant map. A list
with no reset reads one fixed concatenation of fronts and contents.
A reset erases dependence on the input state, so every action with
a reset is constant. If there are $N$ states and at most $K$ actions,
the derived reader has at most $K+N$ actions. For $m$ marker letters,
the junction monitor has $k+m+1$ states and at most $(k+m)m+2$
actions. An empty list gives the identity. A nonempty list with a
failed internal junction gives constant failure. Otherwise the
action tests for one front and sends it to one closing marker,
failing elsewhere. This accounts for all maps on the monitor.

Each word action of the product is a pair of actions. Its number
of states $N'$ and bound $K'$ on actions therefore satisfy
\[
 N'=N(k+m+1),\qquad K'\le (K+N)\bigl((k+m)m+2\bigr).
\]
Since $m\le k$, the sum $K'+N'+1$ is at most
$(K+N+1)(2k^2+2k+4)$. This uses $N'\le N(2k+1)$ and
$(k+m)m+2\le 2k^2+2$. The factor $2k^2+2k+4$ is at most
$F(k)=(2k+3)^3$, by expanding the cube. This also bounds
the growth if a transition table is transported along a
letter bijection, which leaves the word actions unchanged.

\subsection{A dictionary on a reverse trie}

A word can be shortened while preserving its action from every state.
At each step, let $F$ be the map on states of the prefix already
read. Find the longest prefix of what remains that leaves $F$
unchanged and does not consume the last letter. Erase this gap and
retain the next letter. Continue to the end. Each gap is a loop
from every state in the image of its preceding prefix. Hence the
pairs of gap and retained letter give a decomposition with the
same states before each retained letter as in the original run.

The actions of the original prefixes at the beginnings of
different steps are distinct. If a later prefix there had the
same action, the earlier gap could have extended at least as far
as that later prefix, which still excludes the last letter. This
contradicts the choice of the longest gap. There are at most $K$
steps. The last letter is retained, and the argument applies to
any nonempty word. The empty word needs no steps.

Apply this construction on the reverse trie of a finite dictionary
of contents. Mark the root, each dictionary endpoint, and every
vertex with multiple children. Use the paths between consecutive
marks as segments. Every such segment is a word when read toward
the root. Compress it once, valid from every state, and keep its
last letter as in the preceding argument. As read away from the
root, this is the first letter of the shortened segment. Distinct
branches still differ on this first letter. Equality of two
suffixes spelled by retained letters therefore identifies exactly
the same path and the same retained cuts as read from the root.
This means the lists of gaps and retained letters of the two
suffixes coincide. More explicitly, whenever their reverse
paths leave a common marked vertex, different segments have
different first letters. Inside a segment the two paths follow
the same shortened text with the same gap list. Induction along
these paths gives equality of the gap lists up to their shared end
at the retained cut. Thus equal retained suffixes lift to equal
original suffixes at the corresponding cuts.

There are at most $2|\mathcal D|$ segments for a dictionary
$\mathcal D$. If the dictionary is nonempty, each leaf is a
dictionary endpoint, and the number of branching vertices is at
most the number of leaves minus one. The latter follows by summing
outdegrees minus one over the nonleaves of a rooted tree. The root
and the marked endpoints and branch vertices give at most
$1+|\mathcal D|+(|\mathcal D|-1)$ vertices and at most one fewer
segments. An empty dictionary has none. Each compressed segment
has length at most $K$. The retained words have a reverse trie
with at most $2|\mathcal D|K$ edges, since they are read on this
shortened tree. Their lengths are at most this many, and the
ordinary letters in them number at most this many, because each
such letter labels some edge.

Renaming a block by its retained content keeps both markers. This
map is injective on the blocks used, because equality of complete
retained contents lifts to equality of the originals. By validity
of the gaps, the runs before and after retained transitions agree
with their old states. In particular the block action stays the
same, and the closing label agrees. Inside two renamed blocks,
every alignment used by the derived reader lifts within the two
old blocks. Indeed, equality of the content rests identifies
the original gaps, and the label at each matching cut is the
original label at the start or after the preceding retained letter.
If a comparison in the renamed block lists has a nonempty target
rest, injectivity makes the old block rests a comparison as well.
For a final target, the closing label is unchanged, and the host
label again lifts into its block or to its closing label. Consequently
the renaming preserves holes of the derived reader and preserves
linking. It gives expanded words that are subsequences of the
expansions before renaming.

\subsection{The induction for short refinements}

We first shorten holes of arbitrary transition readers. For a depth
bound $d$ reserve
\[
 R\ \ge\ (K+N+1)F(k)^d,\qquad Q=2R+k+2 .
\]
We prove that the shortened words use at most $Q^{d+1}$ distinct
letters in total, and each word has length at most $Q^{(d+1)^2}$.
If every word has length at most one, keep the tuple, which uses
at most $k$ letters. This handles depth zero, where any nonempty
word would itself be a positive rest.

At positive depth, pass to block lists on the product reader.
Their depth is at most $d-1$ if $d\ge 2$. When $d=1$, use the
length-at-most-one case on the block lists. The action and state
bounds allow the same budget $R$ at depth $d-1$ because
$(K'+N'+1)F(k)^{d-1}\le R$. After induction the block lists are
again linked. Shorten the dictionary of contents of their retained
blocks. Its cardinality is at most $Q^d$, because it is the
image of their letter set under taking contents. The transitions
on ordinary letters have at most $K$ actions. On expansion
the support grows to at most $2Q^d K+k\le Q^{d+1}$, and each
block contributes at most $2Q^d K+1\le Q^{d+1}$ letters. These
inequalities use $K\le R$. The lists had length at most
$Q^{d^2}$, so the new words have length at most
$Q^{d^2+d+1}\le Q^{(d+1)^2}$. The block lists were subsequences and
their contents were shortened by deletion. Thus each final word
is a subsequence of its original, and the equivalence on linked
expansions ensures it is in a hole at this level.

To obtain a refinement on a given table, construct a transition
reader from the original tuple. Accept exactly its original
endpoint labels. With a nonempty target rest, admit a pair of labels
if no comparison of that kind in the original tuple between those
coordinates carries them. For a final target, admit any host label
absent from the original host coordinate. The original tuple is a
hole. The new hole keeps the endpoint labels and all its proper
comparisons lift. Every host label must occur in its original
coordinate because it can be compared with a final target, for
example its own. Project labels to states to obtain a refinement.
Summing the length bound and taking $R=(K+N+1)F(k)^d$ proves
\cref{lem:refinement}.

\subsection{The suffix monitor and marked buffers}
\label{app:suffix-monitor}

To keep the terminal windows of a tuple, the suffixes of length at most
$B+2$, let $\mathcal W$ consist of the empty word and the factors of the
terminal windows. This set of window factors is closed under taking
factors. Each window of length $r$ has at most $r^2$ nonempty factors,
by choosing their first and last positions. Thus $|\mathcal W|\le M$,
where $M=1+k(B+2)^2$.

A monitor state is a word $z$ in $\mathcal W$ and a table state
$q$ before $z$. After a letter, the next recorded word is the
longest suffix in $\mathcal W$ of $z$ followed by that letter. This
computes the longest suffix in $\mathcal W$ of the entire input so
far. Indeed, deleting the last letter of any candidate suffix
leaves a suffix in $\mathcal W$ of the previous input, and hence a
suffix of $z$. Advance $q$ by the part of the old recorded word
and new letter before the new suffix. Induction shows that the
record gives the true state before that suffix. It then gives
both the endpoint and the states before all suffixes in $\mathcal W$
at the end of the input.

There are at most $MN$ monitor states. If a word $w$ is in
$\mathcal W$ it contributes at most one action. Otherwise, from any
record $(z,q)$, a suffix in $\mathcal W$ of $zw$ cannot extend left of
$w$, because it would contain $w$ as a factor. So the final recorded
word is the longest suffix in $\mathcal W$ of $w$ alone. The new state
before it is reached by reading $z$ and then the prefix of $w$
before this suffix. The action is determined by one suffix in
$\mathcal W$ and one base action, giving at most $M(K+1)$ actions.
The reasoning applies to every monitor state by regarding $z$
as a word just read from $q$.

In the marked tuple, the fresh terminal marker acts as the identity,
so it does not increase $K$. For every test word the unmarked
stopped suffixes contribute at most $C$ lengths. This follows
by taking the prefix before its first terminal marker, or the
whole test word in the absence of one. Nonempty marked suffixes
contribute at most the one length through the first marker.
The depth is therefore at most $C+1$.

Run the refinement of \cref{lem:refinement} on the suffix monitor
for this marked tuple. Then project its states to those of the
original table by reading the recorded suffix from its stored
state. Equal monitor endpoints keep every original terminal
suffix of length at most $B+2$, with its entry state. For moving
words this keeps the buffer and marker, the last stem letter
when there is one, and the state at the stem's end. If a stem
was empty the entire marked word survives, by the subsequence
condition. Each suffix of length at most $B$ of a shortened
stopped word is also an original suffix, because the original suffix
of that length survives and the length determines it. The total
length bound for the marked refinement bounds each shorter
stopped word and stem by $S$ from \cref{lem:buffered}.
